% Generated draft: requires mathematical and source review.
\documentclass[11pt]{article}
\usepackage[margin=25mm]{geometry}
\usepackage{fontspec}
\setmainfont{Latin Modern Roman}
\usepackage{amsmath,amssymb,mathtools,unicode-math}
\setmathfont{Latin Modern Math}
\usepackage{enumitem,needspace,xurl,hyperref}
\setlength{\parindent}{0pt}
\setlength{\parskip}{5pt}
\setlength{\emergencystretch}{3em}
\newcommand{\N}{\mathbb{N}}
\newcommand{\Z}{\mathbb{Z}}
\newcommand{\Q}{\mathbb{Q}}
\newcommand{\R}{\mathbb{R}}
\newcommand{\C}{\mathbb{C}}
\newcommand{\F}{\mathbb{F}}
\newcommand{\A}{\mathbb{A}}
\newcommand{\PP}{\mathbb{P}}
\newcommand{\E}{\mathbb{E}}
\newcommand{\eps}{\varepsilon}
\newcommand{\dd}{\,\mathrm{d}}
\begin{document}
\section*{UnsolvedMath ID: 30002145}
\textbf{Source problem code:} catalogue entry
\subsection{Definitions and mathematical statement}
For the purpose of explaining the notation, let \(n\geq 1\), let \(\Omega\subseteq\R^n\) be open, and let \(a,b\in\R^n\). These parameters and their hypotheses are not specified in the supplied statement. Define the symmetric tensor product by
\[
 a\mathbin{\odot}b:=\frac{a\otimes b+b\otimes a}{2},
 \qquad (a\otimes b)_{ij}:=a_i b_j,
\]
and the line, possibly degenerate, of symmetric matrices
\[
 L_{a,b}:=\{\lambda(a\mathbin{\odot}b):\lambda\in\R\}.
\]
For a differentiable map \(u:\Omega\to\R^n\), its symmetric gradient is
\[
 Eu:=\frac{Du+(Du)^{\mathsf T}}{2}.
\]
The same expression defines a distribution for locally integrable \(u\). A map belongs to \(\mathrm{BD}(\Omega)\), the space of functions of bounded deformation, when \(u\in L^1(\Omega;\R^n)\) and \(Eu\) is a finite matrix-valued Radon measure.

The supplied problem asks: what structure can be proved for sufficiently regular \(u\) satisfying
\[
 Eu\in L_{a,b}
 =\{\lambda(a\mathbin{\odot}b):\lambda\in\R\}?
\]
Under a function-valued interpretation with fixed \(a,b\), the displayed constraint means that \(Eu(x)\in L_{a,b}\) at the relevant points \(x\in\Omega\); the scalar multiple may depend on \(x\). The source does not specify whether the constraint holds everywhere, almost everywhere, or in a distributional or measure-valued sense. In particular, the reference to BD theory in the background does not specify how this inclusion is to be read when \(Eu\) has a singular measure part.

This is an unresolved formulation of a structure question, rather than a specified theorem to prove. The dimension, domain, regularity class, assumptions and quantifiers on \(a,b\), boundary conditions, and desired structural conclusion are absent. The notation above explains conventional meanings without selecting these missing hypotheses. A precise classification target requires their clarification.
\subsection{Short English statement}
The problem asks what structural restrictions follow when the symmetric gradient of a sufficiently regular vector field lies in the line generated by \(a\mathbin{\odot}b\). The source leaves the regularity, domain, vector hypotheses, meaning of the inclusion, and intended structure unspecified. Its supplied classification as partially solved is dated curation, not an independently verified conclusion.
\subsection{Sources}
\begin{itemize}
\item Supplied local problem-registry record, permanent UnsolvedMath ID \(30002145\), source problem code \texttt{catalogue entry}, title: Rigidity of Symmetric-Gradient Differential Inclusions. Collection: Oberwolfach Reports Open Problems. No individual proposer is supplied; the proposed year is recorded as 2012.
\item Supplied report citation: Calculus of Variations (2013), Oberwolfach Reports, pp.~2205--2268, DOI \texttt{10.4171/owr/2012/36}. \url{https://doi.org/10.4171/owr/2012/36}. The background attributes the problem to workshop open-problem discussions documented in Oberwolfach Reports. The supplied 2013 report label and 2012 proposed year and DOI component are retained without resolving their relationship.
\item Supplied literature source: De Philippis--Rindler, structure of A-free measures. \url{https://arxiv.org/abs/1601.06543}.
\item Supplied literature source: De Philippis--Rindler, symmetric-gradient Young measures. \url{https://arxiv.org/abs/1604.04097}.
\item Supplied literature source: De Philippis--Rindler, fine properties of BD. \url{https://arxiv.org/abs/1911.01356}.
\item Literature assessment dated exactly 2026-08-21: the supplied curation labels the problem partially solved and states that BD structure theorems and good blow-up results provide substantial one-directional rigidity for the symmetric-gradient inclusion, but no complete classification in the generality asked. This is supplied dated curation derived from the OWR record, not independently verified literature assessment; the status and cited claims require independent verification.
\end{itemize}
\end{document}
